% ============================================================================
%  CHAPTER 1: INTRODUCTION   (expanded, deep-dive version)
%  Replaces the entire Chapter 1 body in main.tex (the \chapter line through
%  the end of \section{Novelty}).
% ============================================================================

\chapter{Introduction}
\label{chap:introduction}

Filing an income tax return in India used to be an exercise in copying numbers
from a salary slip into a form. It is no longer that. Over the past decade the
Income Tax Department has assembled a remarkably complete record of each
taxpayer's financial year, and the act of filing has quietly turned into a
problem of making one's own declaration agree with a record the department has
already built. Banks now report interest credited and large deposits. Mutual
funds and stock brokers report every redemption and sale. Crypto exchanges
report transfers, and tax is deducted at source on them under Section 194S.
Employers report salary and the tax they withheld. All of this is collated into
the Annual Information Statement (AIS), a document the department can see in
full but that most taxpayers glance at only after something has gone wrong.

The consequence is a quiet but consequential mismatch. A salaried person sits
down in July, opens a filing portal, types in the figures from Form 16, claims
a handful of deductions, and submits. The software in front of them is competent
at the form and the arithmetic of the slabs, but it is structurally blind to the
AIS. It reasons over what the user typed and nothing else. When the figure the
user declared and the figure the department already holds disagree --- a missed
fixed-deposit interest, an undeclared equity sale, a crypto transfer the filer
forgot about --- the system has no way of noticing, and the first signal the
taxpayer receives is an automated intimation under Section 143(1) some months
later, often with interest and a penalty attached.

A second development has run in parallel. The arrival of capable large language
models has produced a wave of people simply asking a chatbot to work out their
taxes. On the surface this looks like a solution: the model is fluent, it knows
the names of the sections, and it answers in seconds. Underneath, it is unsound
in precisely the place that matters. A general model treats Indian tax
computation as a free-text reasoning task, and it makes confident arithmetic
errors --- a wrong slab boundary, a forgotten cess, a misapplied rebate
threshold --- that a non-specialist has no way of catching. Worse, it will
cheerfully invent a deduction that the law does not permit, such as claiming an
80C investment benefit under a regime that disallows it. The fluency masks the
unreliability, which arguably makes it more dangerous than an honest blank form.

This project, FinITR-AI, grew out of the observation that neither side of this
landscape actually solves the taxpayer's problem. The rule-based tools are
accurate but blind; the chatbots are flexible but wrong. What was missing was a
system that reconciles the three independent views of a person's financial year
--- what the employer reported, what the bank recorded, and what the government
already knows --- before the return is filed, computes the tax with the
exactness of a rule engine, reasons about the unusual cases with the flexibility
of a language model, and refuses to assert anything it cannot ground in a
document. Building that system, and showing that it measurably outperforms the
general models it competes with, is the subject of this report. The remainder of
this chapter states the problem precisely, argues why it is worth solving,
describes the boundary of what we built, lists our concrete objectives, and
explains what is genuinely new in the approach.

\section{Problem Statement}
\label{sec:problem-statement}

At its heart the difficulty is an asymmetry layered on top of another asymmetry.
The first is informational: the taxpayer is asked to self-declare, but is
assessed against a far richer dataset --- the AIS, Form 26AS, and the SFT
(Statement of Financial Transactions) feeds behind them --- that they cannot
easily inspect or interpret. The second is an asymmetry of consequences: an
honest omission and a deliberate evasion look identical to the automated
matching engine, and the cost of either, in interest and penalty and stress,
falls entirely on the individual. The tooling available today does little to
close the first asymmetry, and the fashionable shortcut of asking an AI
introduces a fresh set of failures. Four problems, in particular, motivate this
work.

\begin{itemize}
    \item \textbf{The AIS is ignored at exactly the moment it matters.}
          Conventional filing software and consumer chatbots reason only over
          self-reported figures. Income that the department already knows about
          --- interest, dividends, capital gains, crypto proceeds --- can pass
          silently into an under-declared return. Because the matching happens
          after submission, the taxpayer learns of the gap only when a notice
          arrives, when correcting it is expensive rather than free.
    \item \textbf{Language models cannot be trusted with the arithmetic.}
          Indian tax computation is not arithmetic a model can improvise. It
          involves regime-specific slabs, a standard deduction that differs by
          regime, a Section 87A rebate with a narrow marginal-relief zone, a
          surcharge that itself carries marginal relief, and a four percent cess
          layered on top. General models routinely get at least one of these
          wrong, and they do so with the same confident tone they use when they
          are right.
    \item \textbf{The dual-regime choice is a genuine optimisation, not a
          preference.} Since the New Regime became the default, every filer faces
          a real decision: the Old Regime with its higher rates but full
          deductions, or the New Regime with its lower rates but almost no
          deductions. The right answer depends on the individual's exact
          deduction profile, and an intuition-based guess is frequently the more
          expensive one.
    \item \textbf{Tax documents are among the most sensitive a person owns.}
          Form 16, a full bank statement, and the AIS together paint an
          uncomfortably complete picture of someone's life. Cloud-hosted AI tools
          require uploading precisely this material to a third party, a price
          many users are unwilling --- and, for good reason, ought to be
          reluctant --- to pay.
\end{itemize}

\section{Project Significance}
\label{sec:project-significance}

The significance of FinITR-AI is best understood as attacking all four of those
problems together, rather than improving any one of them in isolation, and doing
so for the ordinary salaried filer rather than only for the tax professional.
Four threads are worth drawing out.

\begin{itemize}
    \item \textbf{It shifts the work from after the notice to before the
          filing.} By reconciling the AIS against Form 16 and the bank statement
          while the return is still being prepared, the system surfaces an
          undeclared item at the only time it can be fixed for free. This turns
          a reactive, defensive process into a preventive one, which is the
          single largest change in the user's experience of filing.
    \item \textbf{The numbers it shows are correct by construction.} Because
          every figure is produced by a deterministic engine that hard-codes the
          current year's rules, the user is not relying on a model having
          happened to memorise the right slab. Correctness stops being a matter
          of probability and becomes a matter of implementation, which is a much
          stronger guarantee for something with financial consequences.
    \item \textbf{Privacy is structural, not promised.} The whole pipeline ---
          parsing, the agents, the local language model served through Ollama,
          and the report generator --- runs on the user's own machine. There is
          no API call carrying a salary figure to a remote server. This is what
          makes the tool usable for the large group of people who would never
          paste their Form 16 into a public chatbot, and it is a property the
          cloud tools simply cannot offer.
    \item \textbf{The output can be audited, not just trusted.} The system does
          not hand back a single number. It produces a reconciliation ledger, a
          risk score with its contributing reasons, a schedule-by-schedule map,
          and a citation-backed explanation --- a brief that a chartered
          accountant, or a careful user, can read line by line and check. An
          answer that can be inspected is worth more, in this domain, than an
          answer that is merely confident.
\end{itemize}

\section{Scope of Work}
\label{sec:scope}

The project is scoped deliberately around the filer who is both the most common
and the most exposed to notices: a salaried individual who may additionally have
income from the capital markets, crypto, interest, or dividends. Within that
boundary the work runs from raw document ingestion all the way to a
filing-ready report, and it is organised around three broad components that the
later chapters develop in full.

\begin{itemize}
    \item \textbf{Ingestion and reconciliation.} A set of parsers reads Form 16
          (as either a PDF or a structured JSON), the AIS, a bank statement in
          CSV form, and optionally Form 26AS, normalising all of them into a
          single shared representation. On top of this sits the reconciliation
          layer that aligns the three views of the year, and a transaction
          classifier that assigns an income type to every line of the bank
          statement.
    \item \textbf{The multi-agent reasoning core.} An orchestrator coordinates
          four specialised agents --- Auditor, Optimizer, Compliance, and Critic
          --- over a shared set of deterministic tools, producing a regime
          recommendation, an ITR-form decision, a schedule map, and a quantified
          notice-risk assessment, with a verification loop that re-runs an agent
          whenever a claim cannot be supported.
    \item \textbf{Output, evaluation, and interface.} A FastAPI backend and a
          Streamlit front-end present the result; an ITR-2 JSON exporter and a
          PDF brief generator produce the deliverables; and a purpose-built
          benchmark, IndianTaxBench, measures the system against GPT-4o and
          Gemini baselines so that the claims made here are quantified rather
          than asserted.
\end{itemize}

Two exclusions keep the project honest about what has actually been built and
tested. FinITR-AI does not file the return on the government portal on the
user's behalf --- it stops at a validated, filing-ready report and an ITR-2 JSON
that mirrors the department's schema. And while it selects ITR-3 and ITR-4 when
a business or presumptive profile is detected, it does not attempt to
comprehensively cover business filers; that lies beyond the salaried-plus-capital
-market scope and is noted as future work rather than claimed as a result.

\section{Objectives}
\label{sec:objectives}

The main objectives of this project are:

\begin{enumerate}
    \item To design a multi-agent architecture that reconciles a taxpayer's
          Form 16, bank statement, and AIS, and that detects the discrepancies
          which would otherwise surface later as a Section 143(1) notice.
    \item To guarantee arithmetic correctness by building a deterministic
          FY 2025-26 tax engine --- covering both regimes, surcharge, cess, the
          Section 87A rebate, and marginal relief --- and to route every
          computation through it rather than through the language model.
    \item To automate the choice between the Old and New regimes and to
          recommend the deductions and CTC-restructuring moves that lawfully
          reduce the user's liability.
    \item To select the applicable ITR form correctly and to map each income
          source to its required schedule using evidence drawn from the
          documents, rather than blanket assumptions.
    \item To suppress hallucinated claims through a dedicated verification agent
          that checks every assertion against the source documents and a
          tax-law corpus, and that forces an upstream agent to re-run when a
          claim cannot be grounded.
    \item To construct a reproducible benchmark, IndianTaxBench, and to use it
          to quantify the system's accuracy and hallucination rate against
          general-purpose language-model baselines.
\end{enumerate}

\section{Novelty}
\label{sec:novelty}

What is new here is not any single piece but the way the pieces are combined to
be accurate, grounded, and private at the same time --- a combination none of
the existing tools or chatbot approaches achieves. Three aspects deserve
particular emphasis.

\begin{itemize}
    \item \textbf{AIS reconciliation is promoted to a first-class step.}
          Instead of treating the AIS as a pre-fill convenience, the Auditor
          agent makes reconciliation between the AIS and Form 16 the very
          entrance to the pipeline, and it converts the outcome into a numeric
          notice-risk score with a per-item breakdown. We are not aware of a
          consumer tool or a prompted model that does this as its primary
          operation.
    \item \textbf{Reasoning and arithmetic are kept strictly apart.} The
          language model is confined to language --- explanation, narrative, and
          fuzzy classification --- while every number is produced by
          deterministic code. This is not a stylistic choice; our ablation
          (Chapter~\ref{chap:results}) shows it is responsible for roughly a
          forty percentage-point difference in tax accuracy, which is the
          difference between a usable system and an unusable one.
    \item \textbf{The verification loop is allowed to say no.} The Critic agent
          does more than attach a confidence score. It checks each claim against
          both the user's documents and a structured tax-law corpus, and it has
          the authority to veto a claim --- an 80C deduction wrongly applied
          under the New Regime, say --- and send the responsible agent back to
          try again. In our measurements this single mechanism cuts the
          hallucination rate from 18.4\% to 6.8\%.
\end{itemize}

Put together, these decisions make FinITR-AI behave less like a chatbot
guessing at taxes and more like a careful junior accountant who reads the source
documents first, refuses to do the slab arithmetic in their head, and has a
senior check every claim before it leaves the building. The chapters that follow
trace how that behaviour is engineered, beginning with a survey of the
approaches it builds on and departs from.
